\section{Dice models}

Dice models are the natural next step after coin models. A die has more than two
visible faces, and this makes it easier to see the difference between a uniform
sample space and a non-uniform probability assignment. Dice also provide one of
the most important examples in elementary probability: the model of two dice,
where elementary outcomes are ordered pairs but many questions are about sums,
comparisons, or other functions of those pairs.

As before, we do not begin with a formula. We begin with the experiment and the
recording rule.

\subsection*{One roll of one die}

Suppose a die is rolled once and the number on the upper face is recorded. The
sample space is
\[
\Omega=\{1,2,3,4,5,6\}.
\]
An elementary outcome is one number from this set.

If the die is modelled as fair, then all six elementary outcomes have the same
probability:
\[
P(1)=P(2)=P(3)=P(4)=P(5)=P(6)=\frac16.
\]
This is the uniform model on a six-element sample space.

If the die is not assumed to be fair, then we may write
\[
P(i)=p_i,\qquad i=1,2,3,4,5,6,
\]
where
\[
p_i\geq 0
\]
for every \(i\), and
\[
p_1+p_2+p_3+p_4+p_5+p_6=1.
\]
The sample space is still \(\{1,2,3,4,5,6\}\), but the probability assignment is
no longer necessarily uniform.

\begin{remarkbox}
The sentence ``a die has six possible results'' does not imply that each result
has probability \(1/6\). The value \(1/6\) comes from the additional modelling
assumption that the die and the rolling mechanism are symmetric enough to treat
the six faces equally.
\end{remarkbox}

\subsection*{Events for one die}

Events are subsets of \(\Omega\). For example, the event that the result is even
is
\[
E=\{2,4,6\}.
\]
The event that the result is at least \(5\) is
\[
A=\{5,6\}.
\]
The event that the result is prime is
\[
B=\{2,3,5\}.
\]

In the fair model,
\[
P(E)=\frac{|E|}{|\Omega|}=\frac36=\frac12,
\]
\[
P(A)=\frac26=\frac13,
\]
and
\[
P(B)=\frac36=\frac12.
\]

In a non-uniform model, we cannot use only the number of favourable outcomes.
Instead, we must add the probabilities of the elementary outcomes inside the
event:
\[
P(E)=p_2+p_4+p_6,
\]
\[
P(A)=p_5+p_6,
\]
\[
P(B)=p_2+p_3+p_5.
\]

\begin{examplebox}
Suppose
\[
(p_1,p_2,p_3,p_4,p_5,p_6)=(0.10,0.10,0.15,0.15,0.20,0.30).
\]
Then the event \(E=\{2,4,6\}\) has probability
\[
P(E)=0.10+0.15+0.30=0.55.
\]
Even though \(E\) contains three of the six possible faces, its probability is
not \(3/6\), because the elementary outcomes are not equally likely.
\end{examplebox}

\subsection*{Two dice as ordered pairs}

Now suppose two dice are rolled and the number on each die is recorded. If the
dice are distinguished, for example by colour, then the outcome is an ordered
pair
\[
(i,j),
\]
where \(i\) is the result on the first die and \(j\) is the result on the second
die. The sample space is
\[
\Omega=\{1,2,3,4,5,6\}\times\{1,2,3,4,5,6\}.
\]
Equivalently,
\[
\Omega=\{(i,j):i,j\in\{1,2,3,4,5,6\}\}.
\]
It has
\[
|\Omega|=6\cdot 6=36
\]
elementary outcomes.

If both dice are fair and the rolls are modelled as independent, then every
ordered pair has probability
\[
\frac1{36}.
\]
Thus the model is uniform on the set of ordered pairs.

\begin{remarkbox}
The outcomes \((2,5)\) and \((5,2)\) are different in this model. They represent
different results on the first and second die. They may lead to the same sum,
but they are not the same elementary outcome.
\end{remarkbox}

\subsection*{The event that the sum is seven}

Consider the event
\[
S_7=\{\text{the sum of the two dice is }7\}.
\]
In the ordered-pair model this event is the set
\[
S_7=\{(1,6),(2,5),(3,4),(4,3),(5,2),(6,1)\}.
\]
It contains six elementary outcomes. Since the fair two-dice model is uniform,
\[
P(S_7)=\frac{|S_7|}{|\Omega|}=\frac6{36}=\frac16.
\]

The count \(6\) comes from a structural fact. To obtain sum \(7\), once the first
coordinate \(i\) is chosen, the second coordinate must be \(7-i\). The possible
first coordinates are \(1,2,3,4,5,6\), and each gives one valid ordered pair.

\begin{examplebox}
The event that the sum is \(4\) is
\[
S_4=\{(1,3),(2,2),(3,1)\}.
\]
Hence
\[
P(S_4)=\frac3{36}=\frac1{12}
\]
in the fair two-dice model.
\end{examplebox}

The events \(S_4\) and \(S_7\) are both described by a sum, but they contain
different numbers of ordered pairs. Therefore they have different probabilities
in the fair two-dice model.

\subsection*{Representing two dice by a table}

A useful representation of the two-dice sample space is a \(6\times 6\) table:
\[
\begin{array}{c|cccccc}
 &1&2&3&4&5&6\\
\hline
1&(1,1)&(1,2)&(1,3)&(1,4)&(1,5)&(1,6)\\
2&(2,1)&(2,2)&(2,3)&(2,4)&(2,5)&(2,6)\\
3&(3,1)&(3,2)&(3,3)&(3,4)&(3,5)&(3,6)\\
4&(4,1)&(4,2)&(4,3)&(4,4)&(4,5)&(4,6)\\
5&(5,1)&(5,2)&(5,3)&(5,4)&(5,5)&(5,6)\\
6&(6,1)&(6,2)&(6,3)&(6,4)&(6,5)&(6,6)
\end{array}
\]
Rows record the first die and columns record the second die. This table is not
just a picture. It makes the elementary outcomes visible and allows events to be
seen as regions.

For example, the event
\[
A=\{(i,j):i=j\}
\]
that both dice show the same number is the diagonal of the table:
\[
A=\{(1,1),(2,2),(3,3),(4,4),(5,5),(6,6)\}.
\]
Thus
\[
P(A)=\frac6{36}=\frac16.
\]

The event
\[
B=\{(i,j):i>j\}
\]
that the first die shows a larger number than the second die is the region below
the diagonal. It contains
\[
1+2+3+4+5=15
\]
elementary outcomes, so
\[
P(B)=\frac{15}{36}=\frac5{12}.
\]

\subsection*{The trap of recording only the sum}

There is another possible recording rule. Instead of recording the ordered pair,
we might record only the sum. Then the sample space becomes
\[
\Omega'=\{2,3,4,5,6,7,8,9,10,11,12\}.
\]
This is a legitimate sample space if the recorded outcome is only the sum.
However, it is not a uniform sample space when the sum comes from rolling two
fair dice.

The value \(7\) can occur in six ordered ways:
\[
(1,6),(2,5),(3,4),(4,3),(5,2),(6,1).
\]
The value \(2\) can occur in only one ordered way:
\[
(1,1).
\]
Therefore
\[
P(7)=\frac6{36},
\qquad
P(2)=\frac1{36}.
\]
They are not equal.

\begin{remarkbox}
The set \(\{2,3,\ldots,12\}\) has eleven elements, but the sums of two fair dice
are not equally likely. A smaller or more aggregated sample space may be useful,
but its probabilities must be inherited from the underlying mechanism, not
assigned uniformly without justification.
\end{remarkbox}

The complete distribution of the sum is
\[
\begin{array}{c|ccccccccccc}
s&2&3&4&5&6&7&8&9&10&11&12\\
\hline
\#\text{ pairs}&1&2&3&4&5&6&5&4&3&2&1
\end{array}
\]
Hence
\[
P(S=s)=\frac{\#\{(i,j):i+j=s\}}{36}.
\]
For example,
\[
P(S=10)=\frac3{36}=\frac1{12}.
\]

\subsection*{Non-uniform dice and product assignments}

If one die has probabilities
\[
P(1)=p_1,\ldots,P(6)=p_6
\]
and another die has probabilities
\[
Q(1)=q_1,\ldots,Q(6)=q_6,
\]
then, under the assumption that the two rolls are independent, the probability
of the ordered pair \((i,j)\) is
\[
P((i,j))=p_iq_j.
\]
The event that the sum is \(7\) then has probability
\[
P(S_7)=p_1q_6+p_2q_5+p_3q_4+p_4q_3+p_5q_2+p_6q_1.
\]
This expression is longer than \(6/36\), but it says exactly the same kind of
thing: add the probabilities of the elementary outcomes that belong to the
event.

\subsection*{Repeated rolls of one die}

For \(n\) rolls of one die, with order recorded, the sample space is
\[
\Omega=\{1,2,3,4,5,6\}^n.
\]
An elementary outcome is a sequence
\[
(a_1,a_2,\ldots,a_n),
\]
where each \(a_i\in\{1,2,3,4,5,6\}\). The number of elementary outcomes is
\[
6^n.
\]
If the die is fair and the rolls are independent, then each sequence has
probability
\[
\frac1{6^n}.
\]
If the die has probabilities \(p_1,\ldots,p_6\), then
\[
P(a_1,a_2,\ldots,a_n)=p_{a_1}p_{a_2}\cdots p_{a_n}.
\]

\begin{examplebox}
For three rolls of a non-uniform die, the sequence \((2,6,2)\) has probability
\[
p_2p_6p_2=p_2^2p_6.
\]
The event ``exactly two sixes in three rolls'' contains the sequences
\[
(6,6,x),\quad (6,x,6),\quad (x,6,6),
\]
where \(x\) is not \(6\). In the fair model its probability is
\[
\binom32\left(\frac16\right)^2\left(\frac56\right).
\]
The binomial coefficient counts the positions occupied by the sixes.
\end{examplebox}

\subsection*{What dice models teach}

Dice models teach three important lessons.

First, a finite sample space may be uniform or non-uniform. The same set
\(\{1,2,3,4,5,6\}\) can describe a fair die or a biased die.

Second, when an experiment has several stages, an elementary outcome is often a
sequence or ordered tuple. For two dice, the natural elementary outcomes are
ordered pairs.

Third, an aggregated value such as a sum is not automatically uniform. The sum
of two fair dice has values in \(\{2,3,\ldots,12\}\), but these values have
different probabilities because they correspond to different numbers of ordered
pairs.

\begin{remarkbox}
The safest modelling habit is to construct the detailed sample space first and
only then decide whether a coarser description, such as the sum, is appropriate.
Aggregation is useful, but it must preserve the probabilities induced by the
underlying model.
\end{remarkbox}
